\documentclass[10pt,a4paper]{amsart}
\usepackage{amsmath,amssymb,amsthm,graphicx}
\usepackage{verbatim}
\newenvironment{reference}{\comment}{\endcomment}
%\newenvironment{reference}{}{}
\newenvironment{slogan}{\comment}{\endcomment}
\newenvironment{history}{\comment}{\endcomment}
% For commutative diagrams we use Xy-pic
\usepackage[all]{xy}

% We use 2cell for 2-commutative diagrams.
\xyoption{2cell}
% Theorem environments.
%
\theoremstyle{plain}
\newtheorem{theorem}[subsection]{Theorem}
\newtheorem{proposition}[subsection]{Proposition}
\newtheorem{lemma}[subsection]{Lemma}

\theoremstyle{definition}
\newtheorem{definition}[subsection]{Definition}
\newtheorem{example}[subsection]{Example}
\newtheorem{exercise}[subsection]{Exercise}
\newtheorem{situation}[subsection]{Situation}

\theoremstyle{remark}
\newtheorem{remark}[subsection]{Remark}
\newtheorem{remarks}[subsection]{Remarks}

\numberwithin{equation}{subsection}

% Macros
%
\def\lim{\mathop{\mathrm{lim}}\nolimits}
\def\colim{\mathop{\mathrm{colim}}\nolimits}
\def\Spec{\mathop{\mathrm{Spec}}}
\def\Hom{\mathop{\mathrm{Hom}}\nolimits}
\def\Ext{\mathop{\mathrm{Ext}}\nolimits}
\def\SheafHom{\mathop{\mathcal{H}\!\mathit{om}}\nolimits}
\def\SheafExt{\mathop{\mathcal{E}\!\mathit{xt}}\nolimits}
\def\Sch{\mathit{Sch}}
\def\Mor{\mathop{\mathrm{Mor}}\nolimits}
\def\Ob{\mathop{\mathrm{Ob}}\nolimits}
\def\Sh{\mathop{\mathit{Sh}}\nolimits}
\def\NL{\mathop{N\!L}\nolimits}
\def\CH{\mathop{\mathrm{CH}}\nolimits}
\def\proetale{{pro\text{-}\acute{e}tale}}
\def\etale{{\acute{e}tale}}
\def\QCoh{\mathit{QCoh}}
\def\Ker{\mathop{\mathrm{Ker}}}
\def\Im{\mathop{\mathrm{Im}}}
\def\Coker{\mathop{\mathrm{Coker}}}
\def\Coim{\mathop{\mathrm{Coim}}}

% Boxtimes
%
\DeclareMathSymbol{\boxtimes}{\mathbin}{AMSa}{"02}

%
% Macros for moduli stacks/spaces
%
\def\QCohstack{\mathcal{QC}\!\mathit{oh}}
\def\Cohstack{\mathcal{C}\!\mathit{oh}}
\def\Spacesstack{\mathcal{S}\!\mathit{paces}}
\def\Quotfunctor{\mathrm{Quot}}
\def\Hilbfunctor{\mathrm{Hilb}}
\def\Curvesstack{\mathcal{C}\!\mathit{urves}}
\def\Polarizedstack{\mathcal{P}\!\mathit{olarized}}
\def\Complexesstack{\mathcal{C}\!\mathit{omplexes}}
% \Pic is the operator that assigns to X its picard group, usage \Pic(X)
% \Picardstack_{X/B} denotes the Picard stack of X over B
% \Picardfunctor_{X/B} denotes the Picard functor of X over B
\def\Pic{\mathop{\mathrm{Pic}}\nolimits}
\def\Picardstack{\mathcal{P}\!\mathit{ic}}
\def\Picardfunctor{\mathrm{Pic}}
\def\Deformationcategory{\mathcal{D}\!\mathit{ef}}
\usepackage{cleveref}
\crefname{originalreference}{original reference}{original references}
\Crefname{originalreference}{Original reference}{Original references}
\makeatletter
\newlabel{categories-lemma-a-version-of-brown}{{categories-lemma-a-version-of-brown}{source}{}{sourcepointer-1}{}}
\newlabel{categories-lemma-a-version-of-brown@cref}{{[originalreference][0][]categories-lemma-a-version-of-brown}{[1][0][]source}{}{}{}}
\makeatother
\crefname{originalreference}{original reference}{original references}
\Crefname{originalreference}{Original reference}{Original references}
\makeatletter
\newlabel{categories-lemma-adjoint-exact}{{categories-lemma-adjoint-exact}{source}{}{sourcepointer-2}{}}
\newlabel{categories-lemma-adjoint-exact@cref}{{[originalreference][0][]categories-lemma-adjoint-exact}{[1][0][]source}{}{}{}}
\makeatother
\crefname{originalreference}{original reference}{original references}
\Crefname{originalreference}{Original reference}{Original references}
\makeatletter
\newlabel{categories-remark-how-to-use-it}{{categories-remark-how-to-use-it}{source}{}{sourcepointer-3}{}}
\newlabel{categories-remark-how-to-use-it@cref}{{[originalreference][0][]categories-remark-how-to-use-it}{[1][0][]source}{}{}{}}
\makeatother
\title{The category of topological groups has colimits, computed on the underlying groups}
\author{The Stacks Project Authors}
\date{stacks-project:topology.tex:local-snapshot-20261012; GNU FDL 1.2 or later}
\begin{document}
\maketitle
\noindent\textit{Editorial note.} All delivered mathematical text is the author's own, copied from the pinned Stacks Project chapter source sources/stacks-topology/pinned-original-source.tex; no statement or proof step was rewritten, reordered or removed.\par\medskip
\noindent\textit{Editorial note.} Preamble provenance: the Stacks Project keeps one single shared preamble for every chapter; this card takes it directly from the staged sources/stacks-topology/preamble.tex as three verbatim spans (lines 8-12, 14-18 and 154-217), so no derived file was written and nothing was added or redefined. The wrapper class and the base packages are supplied by the assembler.\par\medskip
\noindent\textit{Editorial note.} Bibliography provenance: neither the counted statement nor its proof contains a citation (this lane rescans the whole span before writing the card), so the card carries no bibliography and no entry was reconstructed.\par\medskip
\noindent\textit{Editorial note.} Selection provenance: this card was opened under the file-based control contract (CONTROL.md), with the mechanical pre-card gate run on its statement and proof interval before the card was written.\par\medskip
\noindent\textit{Allowed standard prerequisites (editorial).}\begin{itemize}
\item The representability criterion of the form used to construct colimits through a functor commuting with limits. Invoked inside the counted proof (source lines 6113-6133). Produces the representing topological group.
\item The adjoint-exactness statement used to read off commutation with the forgetful functor. Invoked at the end of the counted proof (source lines 6136-6140). Gives the commutation with the forgetful functor to groups.
\end{itemize}\medskip
\ifdefined\hypertarget\hypertarget{sourcepointer-1}{}\fi
\noindent\textit{Source reference.} Original reference categories-lemma-a-version-of-brown: sources/stacks-topology/pinned-original-source.tex, source lines 6099--6099. Navigation pointer only: the counted text refers to the Stacks item categories-lemma-a-version-of-brown, which belongs to another chapter of the project and is not part of this card. The pointer carries the target's own label as its number because this chapter's staged snapshot does not contain the other chapter's tag table; its content is not reproduced here.\par\medskip
\ifdefined\hypertarget\hypertarget{sourcepointer-2}{}\fi
\noindent\textit{Source reference.} Original reference categories-lemma-adjoint-exact: sources/stacks-topology/pinned-original-source.tex, source lines 2585--2585. Navigation pointer only: the counted text refers to the Stacks item categories-lemma-adjoint-exact, which belongs to another chapter of the project and is not part of this card. The pointer carries the target's own label as its number because this chapter's staged snapshot does not contain the other chapter's tag table; its content is not reproduced here.\par\medskip
\ifdefined\hypertarget\hypertarget{sourcepointer-3}{}\fi
\noindent\textit{Source reference.} Original reference categories-remark-how-to-use-it: sources/stacks-topology/pinned-original-source.tex, source lines 5911--5911. Navigation pointer only: the counted text refers to the Stacks item categories-remark-how-to-use-it, which belongs to another chapter of the project and is not part of this card. The pointer carries the target's own label as its number because this chapter's staged snapshot does not contain the other chapter's tag table; its content is not reproduced here.\par\medskip
\begin{lemma}
\label{lemma-topological-group-colimits}
The category of topological groups has colimits and colimits commute
with the forgetful functor to the category of groups.
\end{lemma}

\begin{proof}
We will use the argument of
Categories, Remark \ref{categories-remark-how-to-use-it} to prove
existence of colimits. Namely, suppose
that $\mathcal{I} \to \textit{Top}$, $i \mapsto G_i$ is a functor
into the category $\textit{TopGroup}$ of topological groups.
Then we can consider
$$
F : \textit{TopGroup} \longrightarrow \textit{Sets},\quad
H \longmapsto \lim_\mathcal{I} \Mor_{\textit{TopGroup}}(G_i, H)
$$
This functor commutes with limits. Moreover, given any topological group
$H$ and an element $(\varphi_i : G_i \to H)$ of $F(H)$, there is
a subgroup $H' \subset H$ of cardinality at most $|\coprod G_i|$
(coproduct in the category of groups, i.e., the free product on the $G_i$)
such that the morphisms $\varphi_i$ map into $H'$. Namely, we can
take the induced topology on the subgroup generated by the images
of the $\varphi_i$. Thus it is clear that the hypotheses of
Categories, Lemma \ref{categories-lemma-a-version-of-brown}
are satisfied and we find a topological group $G$
representing the functor $F$, which precisely means that $G$ is
the colimit of the diagram $i \mapsto G_i$.

\medskip\noindent
To see the statement on commutation with the forgetful functor to
groups we will use
Categories, Lemma \ref{categories-lemma-adjoint-exact}.
Indeed, the forgetful functor has a right adjoint, namely the functor which
assigns to a group the corresponding chaotic (or indiscrete) topological group.
\end{proof}

\end{document}
